\documentclass[10pt,a4paper]{article}
\usepackage[fleqn]{amsmath}
\usepackage{amssymb,amsthm,graphicx}
\usepackage{latexsym}
\usepackage[fleqn]{amsmath}
\usepackage{amssymb}
\usepackage{amsthm}
\usepackage{enumitem}
\newcommand{\HCF}{\mathrm{hcf}}
\newcommand{\ZZ}{\mathbb{Z}}
\newcommand{\rk}{\mathrm{rk}}
\DeclareMathOperator{\len}{length}
\DeclareMathOperator{\bl}{b}
\DeclareMathOperator{\cl}{cl}
\DeclareMathOperator{\di}{d}
\DeclareMathOperator{\Cl}{Cl}
\DeclareMathOperator{\acl}{acl}
\DeclareMathOperator{\tp}{tp}
\DeclareMathOperator{\Th}{Th}
\newcommand{\Cor}{\mathcal{C}_\mathrm{or}}
\newcommand{\Cun}{\mathcal{C}_\mathrm{un}}
\newcommand{\mons}{\mathbb{M}}
\newcommand{\strong}{\mathrm{strong}}

\newtheoremstyle{note}% hnamei
{7pt}% hSpace abovei
{10pt}% hSpace belowi
{\normalfont}% hBody fonti
{}% hIndent amounti1
{\bfseries}% hTheorem head fonti
{ }% hPunctuation after theorem headi
{.5em}% hSpace after theorem headi2
{}% hTheorem head spec (can be left empty, meaning ‘normal’)i
\theoremstyle{note}
\newtheorem{thm}{Theorem}[section]
%\newtheorem{theorem}{Theorem}
\newtheorem{lemma}[thm]{Lemma} 
\newtheorem{df}[thm]{Definition}
\newtheorem{rmk}[thm]{Remark} 
\newtheorem{prop}[thm]{Proposition}
\newtheorem{fact}[thm]{Fact}
\newtheorem{cor}[thm]{Corollary}
\newtheorem{conj}[thm]{Conjecture}
\newtheorem{ex}[thm]{Example}
\newtheorem{qu}[thm]{Question}
\newtheorem*{notation}{Notation} 
\title{Repeated-variable reduced substitutions give proper self-embeddings of free Steiner quasigroups}
\author{Silvia Barbina, Enrique Casanovas}
\date{arxiv:2412.06055v3; CC BY 4.0}
\usepackage{url}
\begin{document}
\maketitle
\noindent\textit{Editorial note.} The already public normal-form PropositionMay19\_1 is included as solution support and is not counted again. The new target concerns which substitutions induce proper self-embeddings, a different conclusion.\par\medskip
\noindent\textit{Editorial note.} The final single-occurrence exclusion uses the explicit inverse-term LemmaJune9\_2: if such a term were in the image generated by a,c, its inverse term would put b in that same image. This is the immediate closing implication of the author chain.\par\medskip
\noindent\textit{Editorial scope.} One original proposition with its full author local proof chain; no new claim assigned to any support lemma.\par\medskip
\noindent\textit{Allowed standard prerequisites (editorial).}\begin{itemize}
\item Elementary universal algebra and well-founded recursion. Universal mapping property, generated subalgebras, uniqueness of a homomorphism on generators, and recursion/induction on well-orders.
\item Finite combinatorial accounting on terms and triples. Term length and rank, counting blocks by their greatest vertex, and monotone nonnegative integer sequences.
\end{itemize}\medskip
A \textbf{Steiner triple system} (STS) of order $\lambda$ is a pair $(M, \mathcal{B})$ where $M$ is a set with $\lambda$ elements, and 
 $\mathcal{B}$ is a collection of 3-element subsets of $M$ (the \textbf{blocks})
such that any two $x, y \in M$ are contained in exactly one block. If we require that any two $x,y \in M$ are contained in \textit{at most }one block, $M$ is said to be a  \textbf{partial Steiner triple system}.

From a model theoretic standpoint, Steiner triple systems can be viewed as relational structures with a single a ternary relation $R$ interpreted as the set of blocks (that is, three elements $x, y$ and $z$ are in a block if and only if $R(x,y,z)$). Then $R$ must satisfy
\begin{itemize}
\item for all permutations $\sigma \in S_3$, $ \forall x_1 x_2 x_3 \, \left(R(x_1, x_2, x_3) \rightarrow R(x_{\sigma(1)}, x_{\sigma(2)}, x_{\sigma(3)} \right)$
\item $\forall x y z \left( R(x,y,z) \rightarrow y\neq z \right)$
\item $\forall x y w z \left( R(x,y,w) \wedge R(x,y,z) \rightarrow w=z \right)$.
\end{itemize}
With this choice of language, a substructure of a Steiner triple system is a partial Steiner triple system.

An STS $(M, \mathcal{B})$ determines a \textit{Steiner quasigroup}, that is, a pair $(M, \cdot)$ where $\cdot$ is a binary operation such that
\begin{itemize}
\item $\forall x \, (x\cdot x =x)$
\item if $x \neq y$, then $x \cdot y = z$ where $\{x,y,z \} \in \mathcal{B}$.
\end{itemize}
The operation $\cdot$ defined in this way satisfies
 \begin{enumerate}
\item $\forall x \, (x \cdot y= y\cdot x)$
\item $\forall x \, (x\cdot x =x)$
\item $\forall x \, (x\cdot (x\cdot y)=y)$.
\end{enumerate}
Conversely, any structure $(M, \cdot)$ where $\cdot$ satisfies properties 1, 2 and 3 determines an STS $(M, \mathcal{B})$.

In universal algebraic parlance, Steiner quasigroups are algebras. Thus, if $M$ is a Steiner quasigroup, there are standard notions of a \textit{subalgebra} of $M$, and of the subalgebra \textit{generated} by a subset $A \subseteq M$. A subalgebra of a Steiner quasigroup is a substructure in the model theoretic sense, and in particular it is a Steiner quasigroup. If $A \subseteq M$, then $\langle A \rangle_M$ denotes the substructure of $M$ generated by $A$. The subscript $M$ is usually omitted when clear  from the context.

A Steiner quasigroup induces a partial STS structure on each of its subsets, so we often refer to a partial Steiner triple system when we need to consider a subset of a Steiner quasigroup that is not a substructure.

Given a class  $\mathcal{V}$ of algebras in a given signature, we say that a member $U$ of $\mathcal{V}$ generated by a subset $X \subseteq U$ has the  \textbf{universal mapping property} (UMP) over $X$ with respect to $\mathcal{V}$ if for every $V \in \mathcal{V}$, every map
 \[ \varphi : X \to V \]
extends to a homomorphism $ \hat{\varphi} : U \to V $. A standard argument proves that the homomorphism $\hat{\varphi}$ is unique. Such an algebra $U$ is said to be \textbf{freely generated by $X$}, and $X$ is a \textbf{set of free generators} of $U$. 

It is well known that when $\mathcal{V}$ is a variety, given any cardinal $\lambda \geq 1$ there are $U \in \mathcal{V}$ and $X \subseteq U$ such that $U$ is freely generated by $X $ and $|X|= \lambda$.

Since the class $\mathcal{K}$ of Steiner quasigroups is a variety in the universal algebraic sense, for each cardinal $\lambda \geq 1$, there is a Steiner quasigroup that is freely generated by a set of cardinality $\lambda$. Then a \textbf{free Steiner quasigroup} is a Steiner quasigroup that is freely generated by some subset $X$, and so it has UMP over $X$ with respect to the class~$\mathcal{K}$. We assume $\lambda \geq 3$ to exclude the trivial cases of Steiner quasigroups with one and three elements.

	\begin{df} Let $(M,\cdot)$ be a Steiner quasigroup. A subset $A\subseteq M$ is  \textbf{independent} if the generated substructure $\langle A\rangle_M$ is freely generated from $A$, that is, if $\langle A\rangle_M$ has UMP over $A$ for the class of Steiner quasigroups. 
		
		We say that $A\subseteq M$ is a \textbf{free base} from $M$  if it is independent and it generates $M$. 
	\end{df}
	When we say $\bar{a} = \langle a_i : i \in I \rangle$ is independent, we understand $a_i\neq a_j$ if  $i,j\in I$ and $i\neq j$.
	
	If $M$ is a Steiner quasigroup generated by the subset $A$ and $\bar{a}$ is an enumeration of $A$, then for every $b \in M$ there is a term $t(\bar{x})$ such that $b=t(\bar{a})$.

 Let $M$ be a Steiner quasigroup generated by $A\subseteq M$.  Then  $M=\bigcup_{n<\omega} S_n$ for a chain of sets $S_n$ such that  $S_0 =A$ and $S_{n+1}= \{a\cdot b\mid a,b\in S_n\}$. Note that $S_n\subseteq S_{n+1}$  since  $a\cdot a = a$.
 We will write $S_{<n}$ for $\bigcup_{i<n} S_i$. 
	 The sets $S_n$  are called the \textbf{levels} of $M$ over $A$. An element $a\in M$ has   \textbf{level~$n$ over $A$}  if $a\in S_n\smallsetminus S_{<n}$.

		\begin{df} Let $M$ be a Steiner quasigroup, and let  $A\subseteq M$ be a free base. The sequence $(S_n\mid n<\omega)$ of levels described in Proposition~\ref{May12_2} will be called a \textbf{standard free construction of $M$ over $A$}.
	\end{df}

	\begin{df} We define an equivalence relation $\sim$ on the set of terms in such a way that $t_1 \sim t_2$ if $t_2$ is obtained from $t_1$ by one or more applications of commutativity to subterms.
		
		A term is said to be \textbf{reduced} if it contains no occurrences of subterms of any of the following forms
		\begin{itemize} 
			%\item $t_1t_1$
			%\item $t_1(t_2t_1)$
			\item $t_1 t_2$ with $t_1 \sim t_2$
			\item  $t_1(t_2t_3)$ with $t_1 \sim t_2$ or $t_1 \sim t_3$
			\item $(t_1t_2)t_3$ with $t_1 \sim t_3$ or $t_2 \sim t_3$.
		\end{itemize}
	\end{df}
	
	%The idea is that a reduced term contains no repeated subterms $tt$, not any subterms of the form $t_1(t_2t_1)$.
	
	\begin{df} The \textbf{rank} $\rk(t)$ of a term $t$ is defined to be
		\begin{itemize}
			\item $0$ if $t$ is a variable
			\item $\mathrm{max} \{ \rk (r), \rk(s) \} + 1$ if $t = r\cdot s$.
		\end{itemize}
	\end{df}

\begin{prop}\label{finalprop} Let $\bar{a},b$ be independent, let $r(\bar{x},y)$ be a reduced term with at least two occurrences of $y$, let $c=r(\bar{a},b)$  and let $f:\langle \bar{a},b\rangle\rightarrow \langle\bar{a},c\rangle$ be the surjective homomorphism induced by $f(\bar{a})=\bar{a}$  and $f(b)=c$. Then  $f$ is an isomorphism between $\langle \bar{a},b\rangle$ and $\langle\bar{a},c\rangle$  and  $b\not \in\langle\bar{a},c\rangle$. In fact, for every term $t(\bar{x},y)$ having exactly one occurrence of $y$,  $t(\bar{a},b)\not \in\langle \bar{a},c\rangle$.
\end{prop}

	\begin{prop}\label{May12_2} Let $M$ be a Steiner quasigroup and let $A\subseteq M$. The following are equivalent:
		\begin{enumerate}
			\item $A$ is a free base of $M$.
			\item $M$ is generated by $A$ with  levels $(S_n\mid n<\omega)$ such that $A=S_0$ and
			\begin{enumerate}
				\item if $a,b\in S_n\smallsetminus S_{<n}$ and $a\neq b$, then $a\cdot b\not\in S_n$.
				\item if $a\in S_{n+1}\smallsetminus S_n$, there is a unique pair $\{b,c\}\subseteq S_n$ such that $b\neq c$ and $a= b\cdot c$.
			\end{enumerate}
		\end{enumerate}
		Moreover, property (b) implies that if  $a\in S_{n+1}\smallsetminus S_n$, and $b, a\cdot b\in S_n$, then  $\{b,a\cdot b\}$ is the unique pair given by (b) for  $a$.
	\end{prop}
	\begin{proof} 2 $\Rightarrow$ 1.  Let $N$ be a Steiner quasigroup and let $f:A\rightarrow N$ be a mapping. We extend $f$ to a homomorphism $\hat{f}: M\rightarrow N$.  We obtain $\hat{f}$ as the union $\hat{f}= \bigcup_{n<\omega}f_n$  of a chain  of partial homomorphisms $f_n : S_n \to N$. This means that  $f_n(a\cdot b)= f_n(a)\cdot f_n(b)$  whenever $a,b$ and $a\cdot b$  belong to $\mathrm{dom} (f_n)=S_n$. We start with $f_0=f$.  Recall that $A=S_0$ and that $a\cdot b\not\in S_0$  if $a,b\in S_0$ and, therefore, $f_0$ is a partial homomorphism with domain $S_0$. Now we extend  $f_n:S_n\rightarrow N$  to $f_{n+1}$. %We must define $f_{n+1}(a)$  for every $a\in S_{n+1}\smallsetminus S_n$. 
	For $a \in S_{n+1}\smallsetminus S_n$, there is a unique pair $\{b,c\}\subseteq S_n$ of distinct elements $b,c$ with $a=b\cdot c$. We let  $f_{n+1}(a) = f_n(b)\cdot f_n(c)$.  By (a) and (b), the only defined products involving $a$ in $S_{n+1}$ are   $a=b\cdot c=c\cdot b$, $a\cdot a=a$, $a\cdot b = b\cdot a= c$ and $a\cdot c= c\cdot a =b$ and $f_{n+1}$ respects these products.
		
		1 $\Rightarrow$ 2.  It is easy to construct a Steiner quasigroup $M^\prime\supseteq A $ with the properties described in 2. It is enough to start with $S^\prime_0=A$ and inductively define $S^\prime_{n+1}$ by adding to $S^\prime_n$ a different element $a\cdot b$ for each pair $\{a,b\}$ of elements of $S^\prime_n$ whose product is not defined in $S^\prime_n$. The only  products involving  the new element $c=a\cdot b$ in $S^\prime_{n+1}$ are $c=a\cdot b=b\cdot a$, $c\cdot c=c$, $c\cdot a= a\cdot c=b$ and $c\cdot b=b\cdot c= a$. We have proven that $A$ is a free base of $M^\prime= \bigcup_{n<\omega}S^\prime_n$. The uniqueness of free  Steiner quasigroups over a given set of generators guarantees that $M$ and $M^\prime$ are isomorphic over $A$.  The isomorphism can be used to obtain the levels $S_n$ in $M$ with the required properties.
	\end{proof}

	\begin{lemma}\label{May12_3}\begin{enumerate}
			\item If $t_1\cdot t_2\sim r_1\cdot r_2$, then either  $t_1\sim r_1$ and $t_2\sim r_2$,  or  $t_1\sim r_2$ and $t_2\sim r_1$.
			\item if $t\sim r$ and $t^\prime$ is a subterm of $t$ then $t^\prime\sim r^\prime$ for some subterm $r^\prime$ of $r$.
			\item If $t_1\sim t_2$, then $\rk(t_1)= \rk(t_2)$  and  $t_1(\bar{a})= t_2(\bar{a})$  for any tuple $\bar{a}$.
			\item If  $t_1\sim t_2$ and $t_1$ is reduced, then $t_2$ is reduced.
			\item Any subterm of a reduced term is reduced.
			\item Assume $t,t^\prime$ are reduced, $t\not\sim t^\prime$ and there is no term $r$ such that  $t\sim t^\prime\cdot r$ or $t^\prime\sim t\cdot r$. Then $t\cdot t^\prime$ is reduced.
		\end{enumerate}
	\end{lemma}
	\begin{proof} 1 can be easily checked by induction on the number of applications of commutativity needed for $t_1\cdot t_2\sim r_1\cdot r_2$.  2 is proved by induction on $t$.  3 is clear.  4 follows from 2. 5  and 6 are clear.
	\end{proof}

The next result clarifies how the levels in a construction of a Steiner quasigroup $M$ are linked to the ranks of the reduced terms giving the elements of $M$.
	\begin{lemma}\label{May12_4} Let $M=\langle \bar{a}\rangle$ be a  Steiner quasigroup, with a construction  with levels $(S_n\mid n<\omega)$.
		\begin{enumerate}
			\item If $t$ is a term such that $\rk(t)\leq n$, then $t(\bar{a})\in S_n$.
			\item If $b\in S_n$, then  $b= t(\bar{a})$  for some reduced term $t$ of rank $\rk(t)\leq n$.
			\item  $S_n= \{t(\bar{a})\mid \rk(t)\leq n\} = \{t(\bar{a})\mid t \mbox{ is reduced and } \rk(t)\leq n\}$.
		\end{enumerate}
	\end{lemma}
	\begin{proof} 1 is proved by induction on $n$.  3 follows from 1 and 2. We prove 2 by induction on~$n$.  Assume $n=0$.  If $\bar{a}= \langle a_i : i \in I \rangle$, then $b =a_k$  for some $k \in I$ and $t=x_k$ (a variable, hence reduced and of rank $0$) is the required term. For the inductive step, assume $b\in S_{n+1}$.  If $b\in S_n$ we apply the induction hypothesis. If $b\not\in S_n$, then $b= b_1\cdot b_2$ for $b_1,b_2\in S_n$.  By induction hypothesis, there are reduced terms  $t_1,t_2$ of rank $\leq n$ such that $b_1= t_1(\bar{a})$ and $b_2= t_2(\bar{a})$.  If $t_1\sim t_2$, then $t_1(\bar{a})=t_2(\bar{a})$ and then $b= b_1\cdot b_2 = b_1\cdot b_1 = b_1 \in S_n$, a contradiction. Hence, $t_1\not\sim t_2$.  If there is a term $r$ such that $t_1\sim t_2\cdot r$, then  $b= b_1\cdot b_2=(t_2(\bar{a})\cdot r(\bar{a}))\cdot b_2 = (b_2\cdot r(\bar{a}))\cdot b_2 = r(\bar{a})$. But $\rk(r) <\rk(t_2\cdot r) =\rk(t_1)\leq n$, which implies $b\in S_n$, a contradiction. A similar contradiction follows from the assumption that $t_2\sim t_1\cdot r$ for some term $r$. By  Lemma~\ref{May12_3}, $t= t_1\cdot t_2$ is reduced. Moreover $\rk(t)\leq n+1$  and  $b= t(\bar{a})$.
	\end{proof}

	\begin{prop}\label{May19_1}
		Let $M=\langle \bar{a}\rangle$ be a free Steiner quasigroup with standard free construction $(S_i\mid i<\omega)$, let  $\bar{a}= \langle a_i : i \in I \rangle$ be  independent, and let $t= t(\bar{x})$ be a reduced term of rank $k$. Then
		\begin{enumerate}
			\item $t(\bar{a})\in S_k\smallsetminus S_{<k}$
			\item if $r$ is reduced and  $r(\bar{a})= t(\bar{a})$, then $r\sim t$.
		\end{enumerate}
	\end{prop}
	\begin{proof} We prove 1 and 2 jointly by induction on $k$, but with the extra assumption that $\rk(r)=\rk(t)$. This assumption can be eliminated  later since 1 excludes other possibilities.
		
		\underline{Base case $k=0$}.  If $\rk(t)=0$, then $t$ is a variable $x_i$ and  $t(\bar{a})= a_i\in S_0$.  Moreover, if $r$ has rank 0, $r$ is a variable $x_j$  and  $t(\bar{a})= r(\bar{a})$  implies $a_i =a_j$ and therefore  $x_i = x_j$. Hence  $t\sim r$.
		
		\underline{Inductive step $k+1$}.  Let $t= t_1\cdot t_2$, with $k_1= \rk(t_1)\leq \rk(t_2) = k$ and let  $b= t(\bar{a})$.  By inductive hypothesis, $b_1= t_1(\bar{a})\in S_{k_1}\smallsetminus S_{<k_1}$ and $b_2= t_2(\bar{a})\in S_k\smallsetminus S_{<k}$. Then $b= b_1\cdot b_2$ and there are different cases.
		
		\emph{Case 1}. Assume $k_1 = k$.  If  $b_1 = b_2$, then $b= b_1= b_2$ and by inductive hypothesis $t_1\sim t_2$. But then $t$ is not reduced. Hence, 
		$b_1\neq b_2$. By construction, $b\in S_{k+1}\smallsetminus S_k$.  Moreover, if $r$ is reduced of rank $k+1$ with $r(\bar{a})= b$, then $r= r_1\cdot r_2$, with $r_1,r_2$ reduced and, without loss of generality, $\rk(r_1)\leq \rk(r_2)= k$. By inductive hypothesis $b_1^\prime = r_1(\bar{a})\in S_{\rk(r_1)}\smallsetminus S_{<\rk(r_1)}$  and $b_2^\prime = r_2(\bar{a})\in S_k\smallsetminus S_{<k}$. Then $b= b_1^\prime\cdot b_2^\prime$ and by construction $\{b_1,b_2\}= \{b_1^\prime,b_2^\prime\}$. If $b_1= b_1^\prime$ and $b_2 = b_2^\prime$, then $\rk(r_1)= k_1$ and by inductive hypothesis $t_1\sim r_1$ and $t_2\sim r_2$. Therefore in this case $t\sim r$. If $b_1= b_2^\prime$ and $b_2= b_1^\prime$, then $k=k_1= \rk(r_1)$ and, by inductive hypothesis, $t_1\sim r_2$ and $t_2\sim r_1$. Again we conclude that $t\sim r$.
		
		\emph{Case 2}. Assume $k_1 <k$ and $b\not \in S_{<k}$.  Note that $b\neq b_2$ and $b\not \in S_k$. By construction, $b= b_1\cdot b_2 \in S_{k+1}$.  Now let $r$ be reduced, of rank $k+1$ and such that $r(\bar{a})= t(\bar{a})=b$. Exactly as in the previous case we conclude that $r\sim t$.
		
		\emph{Case 3}. The last case is $k_1 < k$ and  $b\in S_j\smallsetminus S_{<j}$  for some $j<k$.  We show this is impossible.   We know that there is some reduced term $s$ of rank $j$ such that $s(\bar{a}) = b$. Now we claim that $s\cdot t_1$ is reduced. Otherwise, there are two possible cases:
		\begin{enumerate}
			\item[(i)] $s\sim t_1$.  But then $b= b_1$ and hence $b=b_1 = b_2$, which is not the case.
			\item[(ii)]  $t_1\sim s\cdot r$ or $s\sim t_1\cdot r$  for some (reduced)  term $r$.  Then $r(\bar{a})= t_1(\bar{a})\cdot s(\bar{a})= b_1\cdot b= b_2$  and by inductive hypothesis $\rk(r)=k$, contradicting $\rk(t_1)=k_1<k$ and $\rk(s)= j<k$.
		\end{enumerate}
		Therefore, we conclude that $s\cdot t_1$ is reduced. Then, since $\rk(s\cdot t_1)\leq k$ and $s\cdot  t_1 (\bar{a}) = b\cdot b_1 = b_2$, we can apply the inductive hypothesis to get $\rk(s\cdot t_1)= k$  and  $s\cdot t_1\sim t_2$. Then $t\sim t_1\cdot t_2\sim t_1\cdot (s\cdot t_1)$, which says that $t$ is not reduced.
	\end{proof}

	\begin{lemma}\label{June9_2} If $t= t(\bar{x},y)$ has only one occurrence of $y$, then there is a term $r(\bar{x},z)$  such that the equations  $t(\bar{x},y)= z$   and   $r(\bar{x},z)= y$  are equivalent in every Steiner quasigroup.
	\end{lemma}
	\begin{proof} Induction on $t$.  If $t$ is a variable, $t= y$  and we can take  $r= z$.  If $t=t_1\cdot t_2$, without loss of generality, $y$ appears in $t_1$ and not in $t_2$. Then $t(\bar{x},y)=z$ is equivalent to  $t_1(\bar{x},y)= t_2(\bar{x})\cdot z$. By induction hypothesis, there is a term $r(\bar{x},u)$  for which  $t_1(\bar{x},y)= u$  is equivalent to  $r(\bar{x},u)= y$.  Then  $t_1(\bar{x},y)= t_2(\bar{x})\cdot z$  is equivalent to  $r(\bar{x},t_2(\bar{x})\cdot z)= y$.
	\end{proof}
	
	\begin{prop}\label{June9_3} If $t= t(\bar{x},y)$ has only one occurrence of $y$ and $\bar{a}b$ is independent, where $\bar{a}=a_1,\ldots,a_n$, then the mapping induced by
		\begin{itemize}
			\item $a_i\mapsto a_i$  for $i=1,\ldots,n$
			\item $b\mapsto t(\bar{a},b)$
		\end{itemize}
		is an automorphism of  $\langle \bar{a}b\rangle$.
	\end{prop}
	\begin{proof} Since $\bar{a}b$ is independent, there is a homomorphism $f:\langle\bar{a}b\rangle\rightarrow \langle\bar{a}b\rangle$  such that  $f(\bar{a}) = \bar{a}$    and  $f(b)= t(\bar{a},b)$.  We claim that there is another homomorphism $g:\langle\bar{a}b\rangle\rightarrow \langle\bar{a}b\rangle$  such that  $g(\bar{a}) = \bar{a}$   and  $g(t(\bar{a},b))= b$. Then  $g\circ f:\langle\bar{a}b\rangle \rightarrow \langle\bar{a}b\rangle$ is a homomorphism  such that  $g\circ f(\bar{a})= \bar{a}$    and $g\circ f(b)=b$. Hence $g\circ f$ is the identity on $\langle\bar{a}b\rangle$, which implies that $f$ is one-to-one. Lemma~\ref{June9_2}  shows that $b\in\langle \bar{a}, t(\bar{a},b)\rangle$, which is the range of $f$. Hence $f$ is surjective.
		
		We prove the claim by induction on $t$.  If  $t$ is a single variable, it must be  $t= y$ and then $f(b)=b$  and we take $g=f$.  Assume $t= t_1\cdot t_2$. Without loss of generality, $y$ occurs in $t_1$ and not in $t_2$.  The induction hypothesis gives a homomorphism $h:\langle\bar{a}b\rangle\rightarrow\langle \bar{a}b\rangle$ such that  $h(\bar{a})= \bar{a}$ and $h(t_1(\bar{a},b))= b$. Let $a= t_2(\bar{a})$ and note that  $h(a)=a$  and $h(t(\bar{a},b))= h(t_1(\bar{a},b))\cdot h(t_2(\bar{a}))= b\cdot a$. Let $j:\langle\bar{a}b\rangle\rightarrow\langle \bar{a}b\rangle$ be the homomorphism induced by  $j(\bar{a})=  \bar{a}$   and $j(b) =b\cdot a$. Then  $j(b\cdot a)= (b\cdot a)\cdot a =b$ and hence  $j\circ h(t(\bar{a},b))= j(b\cdot a) = b$   and  $j\circ f(\bar{a})= \bar{a}$. Therefore $g= j\circ h$ satisfies the requirements for $t$.
	\end{proof}

\begin{lemma}\label{March_8_23_1} Assume $\bar{x}= \langle x_i : i \in I\rangle$ is a tuple of variables and  $y,z,x_i$ are pairwise distinct. Let  $t_i=t_i(\bar{x},z)$  for $i=1,2$, let $r=r(\bar{x},y)$  and assume $y$ occurs in $r$. We write  $t_i^\prime= t_i(\bar{x},r(\bar{x},y))$  for the substitution, and similarly for other terms.  If  $t_1^\prime\sim t_2^\prime$, then  $t_1\sim t_2$.
\end{lemma}
\begin{proof} By induction on $t_1$.

(i) Let $t_1= x_i$.  Then  $t_1^\prime= x_i\sim t_2^\prime$, which implies $t_2^\prime= x_i$. Since  $y$ occurs in $r$,  $t_2= x_i$ and hence $t_1\sim t_2$.

(ii) Let  $t_1= z$.  Then $t_1^\prime = r\sim t_2^\prime$, which implies that $y$ occurs in $t_2^\prime$ and, therefore, that $z$ occurs in $t_2$.  If $t_2\neq z$, then $\len(t_2^\prime)>\len(r)=\len(t_1^\prime)$, which is impossible if $t_1^\prime\sim t_2^\prime$. Thus, $t_2= z$ and then $t_1\sim t_2$.

(iii) Let  $t_1= s_1\cdot s_2$  for some terms, $s_1,s_2$.  Then $t_1^\prime= s_1^\prime\cdot s_2^\prime\sim t_2^\prime$. Note that $t_2\neq x_i$, since otherwise $t_2^\prime= t_2=x_i\not\sim s_1^\prime\cdot s_2^\prime$. In case $t_2= z$, we get $t_2^\prime=r$, but since $z$ is in  $t_1$ (because $y$ occurs in $t_1^\prime$), then $\len(t_1^\prime)>\len(r)=\len(t_2^\prime)$, contradicting $t_1^\prime\sim t_2^\prime$. Hence, $t_2\neq z$. Therefore, $t_2= u_1\cdot u_2$  for some terms $u_1,u_2$.  Therefore, $t_2^\prime=u_1^\prime\cdot u_2^\prime\sim s_1^\prime\cdot s_2^\prime$. Without loss of generality,  $u_1^\prime\sim s_1^\prime$ and  $u_2^\prime\sim s_2^\prime$. By inductive hypothesis, $u_1\sim s_1$ and $u_2\sim s_2$.  Then $t_1= s_1\cdot s_2\sim u_1\cdot u_2 = t_2$.
\end{proof}

\begin{lemma}\label{March_8_23_2} %Assume $\bar{x}$ is a tuple of variables $x_i$  and  $y,z,x_i$ are all different. 
Assume $\bar{x}= \langle x_i : i \in I\rangle$ is a tuple of variables and  $y,z,x_i$ are pairwise distinct. Let  $t=t(\bar{x},z)$,  let $r=r(\bar{x},y)$  and assume $y$ occurs in $r$. As above, we write  $t^\prime= t(\bar{x},r(\bar{x},y))$  for the substitution. Assume $r=r_1\cdot r_2$ and $y$ occurs in each $r_i$.  If $t$ and $r$ are reduced, then $t^\prime$ is reduced.
\end{lemma}
\begin{proof} By  induction on $t$.

(i) Let $t=x_i$.  Then  $t^\prime = t$  is reduced.

(ii) Let $t=z$. Then $t^\prime = r$ is reduced.

(iii) Let $t= t_1\cdot t_2$  for some terms $t_1,t_2$. By inductive hypothesis, each  $t_i^\prime =t_i(\bar{x},r(\bar{x},y))$  is reduced.  If $t^\prime =t_1^\prime\cdot t_2^\prime$ is not reduced, then either  $t_1^\prime\sim t_2^\prime$  or  $t_1^\prime\sim t_2^\prime\cdot s$  for some term $s$  or  $t_2^\prime\sim t_1^\prime\cdot s$  for some term $s$.  If $t_1^\prime\sim t_2^\prime$, we can apply Lemma~\ref{March_8_23_1} and obtain $t_1\sim t_2$, which implies that $t$ is not reduced.  Now assume, without loss of generality, that $t_1^\prime\sim t_2^\prime\cdot s$  for some term $s$.  We consider three subcases:

\textit{Subcase (a)}: \ $t_1=x_i$. Then  $t_1^\prime = x_i\sim t_2^\prime\cdot s$, which is impossible by length considerations.

\textit{Subcase (b)}: \ $t_1 = z$.  Then $t_1^\prime=r_1\cdot r_2\sim t_2^\prime\cdot s$. The case  $r_1\sim t_2^\prime$  is not possible, since  it implies that $z$ occurs in $t_2$ and then $\len(t_2^\prime)\geq \len (r)>\len(r_1)$. Hence  $r_1\sim s$ and $r_2\sim t_2^\prime$.  But again this is impossible, because it implies that $z$ occurs in $t_2$ and $\len(t_2^\prime)\geq \len(r)>\len(r_2)$.

\textit{Subcase (c)}: \  $t_1= s_1\cdot s_2$  for some terms $s_1,s_2$. Then $t_1^\prime=s_1^\prime\cdot s_2^\prime\sim t_2^\prime\cdot s$. Without loss of generality, $s_1^\prime\sim t_2^\prime$ and $s_2^\prime\sim s$. By  Lemma~\ref{March_8_23_1}, $s_1\sim t_2$. Then $t_1\sim t_2\cdot s_2$  and then  $t\sim (t_2\cdot s_2)\cdot t_2$  is not reduced.
\end{proof}

\begin{proof} Consider first the case  $r= r_1\cdot r_2$, where each $r_i$ contains an occurrence of $y$.  In this case, by Lemma~\ref{March_8_23_2}, for every reduced term $t(\bar{x},z)$,  the term $t^\prime(\bar{x},y)= t(\bar{x},r(\bar{x},y))$ is reduced. 

We prove first that $f$ is one-to-one.  Assume $f(c_1)= f(c_2)$. For each $i$ there is some reduced term $t_i(\bar{x},z)$  such that  $c_i= t_i(\bar{a},b)$.  Since the terms $t_1^\prime,t_2^\prime$ are reduced and  $t_1^\prime(\bar{a},b)= f(c_1)= f(c_2)=t_2^\prime(\bar{a},b)$, by  Proposition~\ref{May19_1},  $t_1^\prime\sim t_2^\prime$.  Then  $t_1\sim t_2$ by Lemma~\ref{March_8_23_1}. Then  $c_1=t_1(\bar{a},b)=t_2(\bar{a},b)= c_2$.

Now we prove that $b\not\in \langle\bar{a},c\rangle$. Assume $b\in\langle\bar{a},c\rangle$. Then there is a reduced term $t(\bar{x},z)$ such that $b= t(\bar{a},c)$. Note that $z$ occurs in $t$, since otherwise $b\in\langle \bar{a}\rangle$.  By Lemma~\ref{March_8_23_2},  $t^\prime(\bar{x},y)= t(\bar{x},r(\bar{x},y))$ is reduced. Clearly, $b= t^\prime(\bar{a},b)$.  If $k=\rk(t^\prime)$, Proposition~\ref{May19_1} gives  $t^\prime(\bar{a},b)\in S_k(\bar{a},b)\smallsetminus S_{<k}(\bar{a},b)$.  But $k\geq \rk(r) >0$ and  $b\in S_0(\bar{a},b)$, a contradiction.

Now we consider the general case of a term $r(\bar{x},y)$  with at least two occurrences of $y$.  We claim that $f$ can be written as a composition $f=g\circ h$, where

(i)  $h$  is an isomorphism between $\langle \bar{a},b\rangle$ and $\langle \bar{a},d\rangle$,  where $d= s(\bar{a},b)$  for a reduced term  $s(\bar{x},y)$  which is a product  $s=s_1\cdot s_2$ of terms $s_i$ containing $y$,   and  $h$ is induced by  $h(\bar{a})= \bar{a}$  and $h(b)=s(\bar{a},b)=d$,  and

(ii) $g$ is an automorphism of  $\langle \bar{a},d\rangle =\langle\bar{a},c\rangle$ and  $c= t(\bar{a},d)$  for a term $t(\bar{x},y)$  with exactly one occurrence of $y$  and  $g$ is induced by $g(\bar{a})=\bar{a}$  and  $g(d)=t(\bar{a},d)= c$.

We prove this by induction on $r$. Clearly, $r$ is not a variable, hence $r=r_1\cdot r_2$  for some reduced terms $r_i= r_i(\bar{x},y)$.  If each $r_i$ contains an occurrence of $y$, we are in the case we have discussed at the beginning and we can take $f=h$ and take $g$ to be the identity on  $\langle \bar{a},c\rangle$. If one of the terms $r_i$, say $r_1$, does not contain $y$, then $r_1= r_1(\bar{x})$  and $r_2(\bar{x},y)$ contains at least two occurrences of $y$. By inductive hypothesis for $r_2$,   we obtain an isomorphism $f^\ast:\langle \bar{a},b\rangle\rightarrow \langle\bar{a},c^\ast\rangle$, induced by  $f^\ast(\bar{a})=\bar{a}$  and $f^\ast(b)=c^\ast= r_2(\bar{a},b)$ which can be decomposed as  $f^\ast = g^\ast\circ h$, where $h:\langle\bar{a},b\rangle \rightarrow \langle\bar{a},d\rangle =\langle \bar{a},c^\ast\rangle$ is induced by $h(\bar{a})=\bar{a}$ and  $h(b)= s(\bar{a},b)=d$, where $s =s_1\cdot s_2$ is reduced and each $s_i $  contains an occurrence of $y$,  and  $g^\ast$ is an automorphism of $\langle\bar{a},d\rangle$ induced by $g^\ast(\bar{a})=\bar{a}$  and  $g^\ast(d)=t^\ast(\bar{a},d)= c^\ast$, where $t^\ast(\bar{x},y)$  contains a unique occurrence of $y$. Now  $t= r_1\cdot t^\ast$  is a term with a unique occurrence of $y$  and  $t(\bar{a},d)= r_1(\bar{a})\cdot t^\ast(\bar{a},d)= r_1(\bar{a})\cdot c^\ast= r_1^\ast(\bar{a})\cdot r_2(\bar{a},b)= r(\bar{a},b)= c$. Hence, there is an automorphism $g$ of $\langle\bar{a},d\rangle$  induced by  $g(\bar{a})=\bar{a}$  and  $g(d)= t(\bar{a},d)= c$  and $f= g\circ h$  is the desired decomposition.
\end{proof}

\begin{thebibliography}{1}

\bibitem{andersenetal}
L.D. Andersen, A.J.W. Hilton, and E.~Mendelsohn.
\newblock Embedding partial {S}teiner triple systems.
\newblock {\em Proceedings of the London Mathematical Society}, 41:557--576,
  1980.

\bibitem{barcas3}
S.~Barbina and E.~Casanovas.
\newblock The elementary theory of free {S}teiner triple systems, 2024.
\newblock \url{http://arxiv.org/abs/2411.13723}.

\bibitem{fujiwara}
T.~Fujiwara.
\newblock Note on the isomorphism problem for free algebraic systems.
\newblock {\em Proc. Japan Acad.}, 32:662--664, 1956.

\bibitem{grishkovetal}
A.~Grishkov, D.~Rasskazova, M.~Rasskazova, and I.~Stuhl.
\newblock Free Steiner triple systems and their automorphism groups.
\newblock {\em Journal of Algebra and Its Applications}, 14(2) 1550025, 2015.

\bibitem{jontarski}
B.~J\'onsson and A.~Tarski.
\newblock On two properties of free algebras.
\newblock {\em Mathematica Scandinavica}, 9(1a):95--101, 1961.

\bibitem{lindner}
C.C. Lindner.
\newblock A survey of embedding theorems for {S}teiner systems.
\newblock {\em Annals of Discrete Mathematics}, 7:175--202, 1980.

\bibitem{sieben}
L.C. Siebenmann.
\newblock A characterization of free projective planes.
\newblock {\em Pacific Journal of Mathematics}, 15(1), 1965.

\bibitem{treash}
C.~Treash.
\newblock The completion of finite incomplete {S}teiner triple systems with
  applications to loop theory.
\newblock {\em Journal of Combinatorial Theory Series A}, 10, 1971.

\end{thebibliography}
\end{document}
